\subsection{1. Introduction}\label{introduction}

A large fraction of machine-learning writing argues that a method
\emph{helps}. A smaller, influential fraction argues that something
\emph{cannot be done}: an information limit, an irreducible error, an
impossibility. Such negative claims are valuable precisely because they
redirect effort: if a task is information-limited, the rational response
is to stop building models and go collect a different signal. They are
also resistant to revision, because they feel like theorems and are
rarely re-examined.

This paper is a fully-worked cautionary study of how a negative claim
can be confidently produced, defended against adversarial scrutiny, and
still be false. It is anchored on one mushroom self-refutation and
replicated across five further look-alike decisions in other kingdoms
(§3.1--§3.2), three of them re-examinations of externally-published
look-alike claims, to show the failure is not domain-specific. We did
not set out to write it. We set out to answer whether ``should I eat
this mushroom?'' is answerable from a photograph, and concluded after
months of work that it is not, that the limit is \emph{information}, not
method, and wrote that paper. Then, preparing it for review, we did the
one thing the original investigation never did: we replaced a
deliberately sovereignty-preserving but weak frozen encoder with a
stronger one. The central result did not survive.

We report the full sequence (the wrong conclusion, the refutation, and
the anatomy of the error) because the failure mode appears general:
safety-negative claims inherit the weakness of their strongest-failing
baseline, and curated benchmarks systematically understate field
difficulty in one direction while overstating model limits in the other.

\textbf{Related work, and our delta.} That weak or poorly-tuned
baselines inflate apparent findings is well documented on the
\emph{positive} side: metric-learning gains that disappear under fair
comparison \citep{musgrave2020}, recommender progress that does not
survive tuned baselines \citep{dacrema2019}, and rank instability across
configurations \citep{dehghani2021, bouthillier2021}. That curated
benchmarks diverge from deployment is the dataset-bias and
distribution-shift program \citep{torralba2011, koh2021, hendrycks2021}.
Our contribution is the \textbf{negative-side analogue and its
asymmetry.} A limitation/impossibility claim inherits the weakness of
its strongest-\emph{failing} baseline exactly as a positive claim
inherits its weakest-\emph{tuned} one. But the community's
instrumentation for the negative case is thinner and concentrated in one
subfield. Adversarial robustness already institutionalizes the move: a
claimed defense or robustness limit is routinely overturned by a
stronger, adaptive attack within the same setup
\citep{athalye2018, tramer2020}. Outside that subfield, and for
frozen-representation limit claims in particular, the practice is
largely absent, and, more sharply, a \emph{self-verification discipline
calibrated to catch rigged wins is structurally blind to over-claimed
losses} (§4d). Two nearest neighbors deserve explicit distinction.
\citet{mcgreivy2024} find that 79\% of ML-for-PDE papers claiming to
beat standard solvers compared against weak baselines, with widespread
under-reporting of negatives. That is the closest positive-side
analogue. Our delta is the \emph{negative-claim} direction and the
representation-invariance instrument that operationalizes it. And on
\emph{this exact task}, \citet{kuznetsov2026} reach the \emph{opposite}
verdict, that AI tools cannot safely identify wild mushrooms (the best
they tested still failed ≈15\% of real-world cases, none consistently
correct), but from deployed, often weak or black-box systems. Our
results re-examine that conclusion and locate the failure in the
deployed model, not an information limit, which is precisely the
distinction this paper exists to draw. We make this concrete with
\textbf{six case studies across five biological groups} (§3--§3.2),
three of them re-examinations of externally-published ``look-alike''
claims, a checklist (§5), and a symmetric ``negative-claim red-team''
protocol (§5.1). We do not claim a measured base rate of false limits in
the literature; establishing one is future work.

\begin{center}\rule{0.5\linewidth}{0.5pt}\end{center}

\subsection{2. The investigation, as we built and believed
it}\label{the-investigation-as-we-built-and-believed-it}

The task. A frozen image encoder over real GBIF photographs of edible
species (chanterelle, field mushroom \emph{Agaricus campestris}, morel)
and poisonous ones (jack-o'-lantern, the Destroying Angel \emph{Amanita
bisporigera}, the open-set Death Cap \emph{A. phalloides}), reduced to a
one-step decision in \{eat, abstain\} with a deliberately asymmetric
cost: eating an edible is a meal, eating a poison is death.

The headline result \emph{as we found it.} On real I-JEPA (ViT-H/14,
mean-pooled) latents, the lethal pair was ``separable'' (AUROC 0.94),
and it did not matter: the actual eat decision cleared \textbf{7 of 81}
held-out Destroying Angels (8.6\% on the split, \textasciitilde4.9\% ±
2.5\% over 20 splits), and \emph{no abstention tuning removed them} (7
eaten at every confidence threshold). The cleared specimens were
confidently-wrong consensus: low predicted risk, low ensemble
disagreement, on-manifold. A 0.94 aggregate AUROC coexisting with a
deadly residual the safety machinery could not catch: the ``94\% → liver
failure'' phenomenon on a real mushroom.

The sweep. We then pushed \textasciitilde18 methods at the residual:
more data, fine-grained metric learning, resampling, the full
f-divergence/optimal-transport family, generative and diffusion
augmentation, metadata fusion, conformal prediction, hierarchical
planning, active-inference exploration, persistent homology, 3D
reconstruction, Morse-theoretic shape analysis, and macro↔micro
generative modeling. Each was run with fair baselines; the most
important claims also with adversarial verification. None removed the
residual. Decision-curve analysis showed that at the field-standard
lethal cost ratio (FungiCLEF \texttt{c\_PSC\ =\ 100}) no operating point
beat ``eat nothing.'' A frontier VLM refused every deadly specimen, but,
we argued, from world knowledge of famous species, not from overcoming
the underlying limit.

The theory. We formalized the negative (Appendix A): a
\emph{coincidence/divergence floor} (where two class-conditional feature
distributions coincide, every f-divergence and Wasserstein vanish and
Bayes error is the prior), a \emph{data-processing bound} (any function
of the photo, whether a sharper representation, a generated image, an
upscale, or a topological summary, is bounded by the photo's Bayes
floor), and an \emph{independence condition} (only a channel with
information about edibility \emph{beyond} the photo can help). These
say: if the discriminating feature is not in the image, no method
recovers it.

The discipline. Twice during this work we produced a result that
\emph{looked} like a method win and was not: a hierarchical-planning
improvement that adversarial-verification panels decomposed as a rigged
baseline and an objective swap. We caught both, reported them, and
adopted adversarial verification as our standing practice. We were, we
believed, being unusually careful.

\textbf{This is the paper we were ready to submit.} It is reproduced,
essentially intact, as the companion document; we preserve it
deliberately as the specimen.

\begin{center}\rule{0.5\linewidth}{0.5pt}\end{center}

\subsection{3. The refutation}\label{the-refutation}

Preparing for review, a referee-style objection was obvious: \emph{the
backbone is a frozen, mean-pooled I-JEPA; did you try a stronger
representation?} We had, in fact, shown that fine-tuning a domain
encoder helped \emph{in-category}, but we had filed that under ``ceiling
up, floor stays'' and not pursued frozen alternatives. So we ran the
sweep we should have run first.

\textbf{Five real pairs × three strong frozen encoders.} We assembled
five genuine field-confusion edible/deadly pairs
(field-mushroom/Destroying-Angel, chanterelle/jack-o'-lantern,
morel/false-morel, oyster/jack-o'- lantern, puffball/death-cap-button)
and encoded each with three paradigmatically different strong frozen
encoders: masked-latent \textbf{I-JEPA} \citep{assran2023} (ViT-H/14),
self-distilled \textbf{DINOv2} \citep{oquab2024} (ViT-L/14), and
contrastive \textbf{CLIP} \citep{radford2021} (ViT-L/14). Per cell, the
actual false-eat at 80\% edible-yield, 10 seeds.

\textbf{Decision rule.} On each encoder's frozen embeddings we fit the
identical probe: an ensemble of small multilayer-perceptron risk nets
(three linear layers with 128 and 32 hidden units and a sigmoid output
giving the probability of ``poisonous''), each trained on the per-seed
30\% train split. The ensemble mean is the risk score and its standard
deviation is the disagreement referred to in §2. The eat/abstain rule
thresholds the ensemble-mean risk at the operating point that gives 80\%
edible yield on the split. Because the same probe is applied to every
encoder, a within-pair difference in false-eat is attributable to the
representation and not to the read-out. Exact encoder checkpoints and
probe hyperparameters are in Appendix C.

{\def\LTcaptype{none} % do not increment counter
\begin{longtable}[]{@{}
  >{\raggedright\arraybackslash}p{(\linewidth - 6\tabcolsep) * \real{0.2500}}
  >{\raggedright\arraybackslash}p{(\linewidth - 6\tabcolsep) * \real{0.2500}}
  >{\raggedright\arraybackslash}p{(\linewidth - 6\tabcolsep) * \real{0.2500}}
  >{\raggedright\arraybackslash}p{(\linewidth - 6\tabcolsep) * \real{0.2500}}@{}}
\toprule\noalign{}
\begin{minipage}[b]{\linewidth}\raggedright
mean over 5 pairs
\end{minipage} & \begin{minipage}[b]{\linewidth}\raggedright
I-JEPA-H14
\end{minipage} & \begin{minipage}[b]{\linewidth}\raggedright
DINOv2-L
\end{minipage} & \begin{minipage}[b]{\linewidth}\raggedright
CLIP-L
\end{minipage} \\
\midrule\noalign{}
\endhead
\bottomrule\noalign{}
\endlastfoot
false-eat at 80\% & 1.3\% (95\% UCB 2.2\%, n=720) & 0.1\% (95\% UCB
0.4\%, n=720) & 0.0\% (95\% UCB 0.4\%, n=720) \\
open-set eaten & 27\% & 0\% & 3\% \\
\end{longtable}
}

\emph{Bounds are one-sided 95\% Clopper-Pearson on the false-eat rate.
\texttt{n=720} is the five pairs \textbf{pooled}, with
\textasciitilde150 unique held-out poison specimens per pair (the
per-seed 30\% test split over }all available images\emph{, \textbf{not}
seed-multiplied; the 10 seeds reshuffle the same pool). The strong
encoders are \textbf{near-zero but not exactly zero}: DINOv2's pooled
estimate is 0.1\% (its Agaricus/Destroying-Angel cell alone reads
0.3\%), CLIP's 0.0\%, both with pooled upper bound ≈0.4\%. At the
original ≈60-specimen-per-cell sampling these same cells read a bare
``0.0\%'' with a ≈4.9\% per-cell upper bound that }overlapped* the
original \textasciitilde4.9\% ± 2.5\% residual (§2). That is, the exact
zero was itself an undersampling artifact, the very effect this paper
warns against (§4). The encoder gap is unchanged: weak I-JEPA's pooled
false-eat is 1.3\% (UCB 2.2\%), an order of magnitude above the strong
encoders.*

The residual did not persist across encoders. \textbf{DINOv2 and CLIP
drove the in-distribution false-eat to near zero: DINOv2 ≤0.3\% on every
one of the five pairs (pooled 0.1\%, UCB 0.4\%), CLIP ≤0.1\%.} On the
\emph{identical} Agaricus-vs-Destroying-Angel task where I-JEPA leaves
3.5\%, DINOv2 leaves 0.3\% (UCB 2.0\%). A center-crop confound test
(background removed) left DINOv2 essentially unchanged (≤0.2\% per cell,
pooled 95\% UCB 0.4\%), so it was reading the fungus, not the habitat.

\textbf{Field degradation.} We tested whether the curated success
survives the regime where foraging deaths occur: downsample, blur,
\textbf{occlude the bottom third} (the un-excavated base, the single
most common field failure, which removes the volva), and color-jitter
the genuine lookalike pairs. DINOv2 held at \textbf{0.0--0.7\%} (95\%
UCB ≤3.1\%, n=150/cell); a lookalike open-set of novel white toxins
(n=1151) was eaten 0\% (95\% UCB 0.5\%).

\textbf{The decisive test (the hardest pairs we could find).} The above
pairs are separable when photographed well. A fair test of an
\emph{information} limit needs the cases mycologists call genuinely
undecidable from a photo. We freshly scraped the textbook
``requires-a-microscope'' confusions and tested them with the easy pair
as a positive control:

{\def\LTcaptype{none} % do not increment counter
\begin{longtable}[]{@{}
  >{\raggedright\arraybackslash}p{(\linewidth - 6\tabcolsep) * \real{0.2500}}
  >{\raggedright\arraybackslash}p{(\linewidth - 6\tabcolsep) * \real{0.2500}}
  >{\raggedright\arraybackslash}p{(\linewidth - 6\tabcolsep) * \real{0.2500}}
  >{\raggedright\arraybackslash}p{(\linewidth - 6\tabcolsep) * \real{0.2500}}@{}}
\toprule\noalign{}
\begin{minipage}[b]{\linewidth}\raggedright
pair (10-seed AUROC / false-eat at 80)
\end{minipage} & \begin{minipage}[b]{\linewidth}\raggedright
I-JEPA
\end{minipage} & \begin{minipage}[b]{\linewidth}\raggedright
DINOv2
\end{minipage} & \begin{minipage}[b]{\linewidth}\raggedright
CLIP
\end{minipage} \\
\midrule\noalign{}
\endhead
\bottomrule\noalign{}
\endlastfoot
{[}control{]} Agaricus / Destroying-Angel & 0.97 / 3.5\% & 1.00 / 0.3\%
(95\% UCB 2.0\%, n=150) & 1.00 / 0.1\% (95\% UCB 2.0\%, n=150) \\
\textbf{honey-mushroom / deadly \emph{Galerina}} (most fatal real
confusion) & 0.97 / 3.5\% & \textbf{1.00 / 0.0\% (95\% UCB 2.0\%,
n=150)} & 1.00 / 0.1\% (95\% UCB 2.0\%, n=150) \\
\textbf{woodtuft / deadly \emph{Galerina}} (``separable only by spore
ornamentation'') & 0.95 / 4.8\% & \textbf{0.98 / 0.7\% (95\% UCB 3.1\%,
n=150)} & 0.98 / 1.9\% (95\% UCB 5.1\%, n=150) \\
honey-mushroom / sulphur-tuft & 0.88 / 21.3\% & 1.00 / 0.4\% (95\% UCB
2.5\%, n=120) & 0.98 / 1.4\% (95\% UCB 5.2\%, n=120) \\
puffball / earthball & 0.98 / 2.1\% & 1.00 / 0.0\% (95\% UCB 2.5\%,
n=120) & 1.00 / 0.4\% (95\% UCB 2.5\%, n=120) \\
\end{longtable}
}

\emph{Bounds are one-sided 95\% Clopper-Pearson on the false-eat rate;
\texttt{n} is the number of \textbf{unique held-out poison specimens}
(the per-seed 30\% test split over }all available images\emph{,
\textbf{not} seed-multiplied). \texttt{n} differs by row because the
species differ in image count: the control and both }Galerina* rows have
\texttt{n=150}, the sulphur-tuft and earthball rows \texttt{n=120}. Two
cells (honey/\emph{Galerina} and puffball/earthball, DINOv2) hold at an
exact 0 of 150/120; the rest are \textbf{small but nonzero} (DINOv2
0.3--0.7\%, CLIP 0.1--1.9\%). At the original ≈60--90-specimen sampling
several of these read a bare ``0.0\%'', so the higher-powered re-run
shows those zeros were undersampling artifacts. The weak I-JEPA
baseline's one-sided 95\% upper bounds run 6.3--28.8\%. Point estimates
are the full-image re-run; the encoder-dependence conclusion is
unchanged.*

DINOv2 performed nearly perfectly on the pairs we selected
\emph{because} they are supposed to be undecidable from a photo. A
substrate-stripping crop (the one plausible reprieve, since
\emph{Armillaria} on living hardwood vs \emph{Galerina} on conifer logs
is a habitat cue) left DINOv2 low (\textbf{0.0--1.0\%}, 95\% UCB
≤4.1\%), while the \emph{weak} I-JEPA stayed high and degraded further
on the canonical pair (honey/\emph{Galerina} 3.5→4.1\%). This confirmed
that the weak encoder relied on contextual cues that the strong encoder
did not require. \emph{(All sub-tests above, namely field degradation,
±background crop, and substrate-strip crop, were re-run on all available
images, matching the two main pair tables. DINOv2's false-eat stays
≤1.0\% per cell across every regime, CLIP ≤3.8\% (only under the
substrate-strip crop).)}

\textbf{Across three regimes (curated pairs (± background crop), field
degradation with base occlusion, and open-set) spanning nine lethal
pairs, all re-run on every available image, the strong encoder held the
false-eat to ≤1.0\% (DINOv2; ≤0.7\% un-cropped, curated-pooled 95\% UCB
0.4\%) and ≤3.8\% (CLIP, under substrate-strip crop) where the weak
encoder left 0.1--21.3\%, lower within every pair.} The curated GBIF
pairs are the benchmark regime; the field-degraded, base-occluded, and
open-set results are the harder regime where deaths occur, and the
worst-case individual specimen (§6) remains untested. The decisive logic
is not the absolute rate (the strong-encoder cells are near-zero, not
exact zeros) but its \emph{encoder-dependence}: an information limit
binds the photograph, hence every encoder equally, so DINOv2 beating
I-JEPA \emph{within every pair} is incompatible with one. The
\textasciitilde5\% ``information-limited residual'' the original paper
was built on tracked the backbone, not the photo.

\subsubsection{3.1 A second case, to test whether the failure is
mushroom-specific}\label{a-second-case-to-test-whether-the-failure-is-mushroom-specific}

A single self-refutation invites the obvious objection: perhaps
mushrooms are special, or our weak encoder was uniquely ill-suited to
them. So we ran the identical protocol on a \emph{second, independent
lethal decision in a different kingdom}: the deadliest look-alike
confusions in the carrot family (Apiaceae). These are \textbf{poison
hemlock} (\emph{Conium maculatum}, the plant that killed Socrates) and
\textbf{water hemlock} (\emph{Cicuta maculata}, the most violently toxic
plant in North America) against edible \textbf{wild carrot / Queen
Anne's lace} (\emph{Daucus carota}) and \textbf{cow parsley}
(\emph{Anthriscus sylvestris}), a confusion that kills foragers every
year. Fresh GBIF scrape, same three encoders, same eval, same
Clopper-Pearson bounds (n = 90 held-out poison specimens per pair):

{\def\LTcaptype{none} % do not increment counter
\begin{longtable}[]{@{}
  >{\raggedright\arraybackslash}p{(\linewidth - 6\tabcolsep) * \real{0.2500}}
  >{\raggedright\arraybackslash}p{(\linewidth - 6\tabcolsep) * \real{0.2500}}
  >{\raggedright\arraybackslash}p{(\linewidth - 6\tabcolsep) * \real{0.2500}}
  >{\raggedright\arraybackslash}p{(\linewidth - 6\tabcolsep) * \real{0.2500}}@{}}
\toprule\noalign{}
\begin{minipage}[b]{\linewidth}\raggedright
pair (10-seed AUROC / false-eat at 80, 95\% UCB)
\end{minipage} & \begin{minipage}[b]{\linewidth}\raggedright
I-JEPA
\end{minipage} & \begin{minipage}[b]{\linewidth}\raggedright
DINOv2
\end{minipage} & \begin{minipage}[b]{\linewidth}\raggedright
CLIP
\end{minipage} \\
\midrule\noalign{}
\endhead
\bottomrule\noalign{}
\endlastfoot
poison-hemlock / wild-carrot & 0.93 / \textbf{9.0\%} (≤15.5\%) & 1.00 /
0.2\% (≤3.3\%) & 0.99 / 0.2\% (≤3.3\%) \\
water-hemlock / wild-carrot & 0.96 / \textbf{5.8\%} (≤11.3\%) & 1.00 /
0.0\% (≤3.3\%) & 1.00 / 0.2\% (≤3.3\%) \\
poison-hemlock / cow-parsley & 0.90 / \textbf{16.2\%} (≤24.5\%) & 0.99 /
0.8\% (≤5.2\%) & 0.98 / 1.4\% (≤5.2\%) \\
\end{longtable}
}

The pattern reproduces, more clearly than in mushrooms. The weak encoder
leaves a \textbf{measured 5.8--16.2\% false-eat} (a residual that,
reported alone, would read as a substantial information limit on a
deadly decision), and a stronger frozen encoder reduces it to 0--0.8\%
within every pair (the I-JEPA one-sided 95\% lower bounds, 2.2--10.6\%,
exceed the DINOv2 point estimates in every pair, and clear the DINOv2
upper bounds in two of the three pairs, with water-hemlock / wild-carrot
the lone overlap: I-JEPA lower 2.2\% against DINOv2 upper 3.3\%). The
representation-invariance test (§5.1, item 2) flags both domains
identically: a limit one encoder sees and another does not is a property
of the encoder, not the photograph. The failure is \textbf{not
mushroom-specific}; it is a generic hazard of anchoring a
safety-negative on a weak representation. (Both these cases are our own
constructions; §3.2 takes the next step and re-examines three
\emph{externally published} indistinguishability claims.)

\subsubsection{3.2 The same protocol across five biological groups:
three external
claims}\label{the-same-protocol-across-five-biological-groups-three-external-claims}

To test the failure mode at scale, we ran the identical three-encoder
protocol on a battery of well-known lethal look-alike confusions
spanning five biological groups, \textbf{three of them anchored to
externally-published claims} (two of visual indistinguishability, one of
systematic human misidentification), so the ``limit'' is not one we
authored:

{\def\LTcaptype{none} % do not increment counter
\begin{longtable}[]{@{}
  >{\raggedright\arraybackslash}p{(\linewidth - 6\tabcolsep) * \real{0.2500}}
  >{\raggedright\arraybackslash}p{(\linewidth - 6\tabcolsep) * \real{0.2500}}
  >{\raggedright\arraybackslash}p{(\linewidth - 6\tabcolsep) * \real{0.2500}}
  >{\raggedright\arraybackslash}p{(\linewidth - 6\tabcolsep) * \real{0.2500}}@{}}
\toprule\noalign{}
\begin{minipage}[b]{\linewidth}\raggedright
case (group)
\end{minipage} & \begin{minipage}[b]{\linewidth}\raggedright
external ``look-alike'' claim
\end{minipage} & \begin{minipage}[b]{\linewidth}\raggedright
I-JEPA false-safe at 80\%
\end{minipage} & \begin{minipage}[b]{\linewidth}\raggedright
DINOv2
\end{minipage} \\
\midrule\noalign{}
\endhead
\bottomrule\noalign{}
\endlastfoot
mushrooms (fungi) & none (this work) & 0.1--21.3\% & 0--0.7\% \\
hemlock / wild carrot (plants) & none (this work) & 5.8--16.2\% &
0--0.8\% \\
coral snake / kingsnake (reptiles) & \citep{pfennig2001} & 4.2--22.4\% &
\textbf{0.0\%} \\
brown recluse / house spider (arachnids) & \citep{vetter2000} &
10.3--\textbf{48.9\%} & 0--2.0\% \\
poison-frog mimic / model (amphibians) & \citep{symula2001} &
12.7--\textbf{43.3\%} & 0.3--10.0\% \\
deadly berries (plant fruit) & none (foraging literature) & 0.2--19.9\%
& \textbf{0.0\%} \\
\end{longtable}
}

\emph{(Full per-pair tables with Clopper-Pearson bounds in Appendix C.
Per our own no-bare-zeros rule (§5.1, item 5), the DINOv2 ``0.0\%''
entries are 0 observed of ≈90 held-out poison specimens per pair. Each
is a one-sided 95\% upper bound ≤3.3\%, not exact zeros. The amphibian
row is small-n, see below.)} The pattern holds across all six groups:
the weak encoder leaves a large, measured false-safe rate, frequently
\textbf{near-chance} (the recluse / house-spider and the
\emph{fantastica} / \emph{imitator} pairs reach AUROC 0.73 and 0.72,
clearing roughly half of the dangerous specimens), and a stronger frozen
encoder lowers it within every pair, to near-zero in every group but the
small-sample poison frogs (qualifier iii). A reviewer shown only the
I-JEPA column would reasonably conclude that several of these decisions
sit at an information limit. For the three external rows the published
literature broadly agrees, in two distinct senses. For the coral snakes
and poison frogs it asserts visual near-identity: the coral snake and
its kingsnake mimics are a canonical Batesian look-alike pair, their
resemblance the premise of the frequency-dependent mimicry Pfennig et
al.~2001 established, and poison-frog mimics co-evolved to be visually
identical to their models (Symula 2001). The brown recluse is a
different kind of anchor: Vetter documents that recluses are massively
over-reported because non-experts rely on the unreliable violin marking
rather than the diagnostic eye arrangement, a systematic
human-misidentification rather than a claim of visual
indistinguishability. A strong encoder separating the recluse pair is
therefore consistent with Vetter, while it violates the
visual-indistinguishability premise the snakes and frogs are supposed to
satisfy. In every row the weak encoder reproduces the human-level
confusion and the strong encoder does not.

Three honest qualifiers. (i) These are species distributions on (often
whole-organism) GBIF photographs, so they test \emph{species
separability}, not the isolated diagnostic confusion a forager faces
with a single berry cluster or a partial specimen. The field and
worst-case caveats of §6 apply unchanged. (ii) The pattern is
\textbf{not universal} (the bilberry / deadly-nightshade pair is easy
even for the weak encoder (0.2\%)), and that is the point: the
encoder-dependence appears precisely where a weak baseline \emph{would}
have been read as a limit, and is absent where it would not. The
representation-invariance test (§5.1, item 2) separates the two
automatically. (iii) The poison frogs are the one group where the
strong-encoder advantage, though present in every pair point-wise, is
\textbf{not yet statistically decisive}: the mimic/model species are
photograph-poor, giving only n=22--30 held-out specimens on three of the
four pairs (DINOv2 up to 10.0\%, upper bound ≤25.9\%), so only the n=90
\emph{Allobates} / \emph{Ameerega bilinguis} pair is conclusive. We
report the frog case as suggestive, not confirmatory (per-pair bounds in
Appendix C.1).

\begin{center}\rule{0.5\linewidth}{0.5pt}\end{center}

\subsection{4. Anatomy of the error}\label{anatomy-of-the-error}

How did a careful, reproduced, adversarially-verified investigation
produce and defend a false limit? Four compounding causes, each general.

\textbf{(a) The baseline was weak, and a negative inherits its
baseline.} A claim that ``no method beats error ε'' is exactly as strong
as the \emph{best} method that fails to. We used frozen mean-pooled
I-JEPA for principled reasons (sovereignty, a fixed feature space for
the divergence-floor proof) and never asked whether it was \emph{strong
enough to anchor a negative}. It was not. Under a stronger frozen
encoder the same residual was an order of magnitude smaller (near zero,
≤0.7\% per cell): it was a property of the backbone, not the photograph.
An ablation (\texttt{ablation\_agg.py}) isolates the cause, and finds
that the baseline was weak in \emph{two} ways at once. First, the
pretraining objective: replacing I-JEPA's mean-pool with max-pool does
not rescue it (both stay near-chance, 51\% and 56\% false-safe on the
recluse pair), so I-JEPA's masked-latent features are weak for this
discrimination regardless of how they are pooled. Second, the
aggregation is independently lossy: mean-pooling DINOv2's patch tokens
rather than using its CLS embedding degrades that same hard pair from
2.1\% to 25.7\% (max-pool recovers most, 5.5\%). Our original baseline
combined both handicaps, and we checked neither.

\textbf{(b) Curated data was mistaken for the field.} GBIF
research-grade photographs are \emph{selected to show the diagnostic
features}. A model that scores \textasciitilde0 on them has learned that
the species differ, which they do, in photographable ways, even for
``microscope'' pairs (the microscope is needed for \emph{certain}
identification, not for a good \emph{probabilistic} guess). We treated
benchmark separability as a proxy for field decidability. It runs the
other way: curated benchmarks understate field difficulty and overstate
the model's limits at once.

\textbf{(c) The theory was true but empirically vacuous.} Our formal
floor and data-processing propositions are correct: \emph{if} the photo
lacks the discriminating information, no function of it recovers it. We
never verified the antecedent. For every real lethal pair we could
obtain, the antecedent is false at the distribution level: the photo
\emph{does} carry enough signal, so the propositions, while valid, said
nothing about these pairs. A true theorem with an unverified antecedent
is not evidence.

\textbf{(d) Adversarial verification checked the wrong axis.} Our
verification panels were genuinely effective: they reproduced every
number and caught two rigged \emph{positive} results. But they audited
\emph{correctness} (do the numbers reproduce?) and \emph{self-deception
toward positives} (is this win rigged?). They never asked the question
that mattered for a \emph{negative}: \emph{is the baseline strong enough
that its failure is informative?} A verification discipline tuned to
catch over-claimed wins is blind to over-claimed losses.

The unifying point is that we instrumented heavily against finding a
false positive and not at all against asserting a false negative. That
asymmetry is the central lesson. We establish it from six case studies
across five biological groups (three re-examining externally-published
look-alike claims). This is an existence proof and a protocol, not a
measured base rate over the literature; how often published limit claims
are refuted by the strong-baseline test is the question §5.1
operationalizes.

\begin{center}\rule{0.5\linewidth}{0.5pt}\end{center}

\subsection{5. A checklist for safety-negative
claims}\label{a-checklist-for-safety-negative-claims}

From the post-mortem, the stress tests a limitation/impossibility claim
should pass before it is trusted:

\begin{enumerate}
\def\labelenumi{\arabic{enumi}.}
\tightlist
\item
  \textbf{Strong-baseline floor.} Re-run the negative against the
  \emph{strongest reasonable} baseline, not the most principled or
  convenient one. A residual that vanishes under a stronger backbone was
  never a limit. Report the strongest baseline that \emph{fails}, and
  treat the claim as bounded by it.
\item
  \textbf{Antecedent verification.} If the negative rests on a theorem
  (``no function of X helps''), test that the theorem's antecedent (``X
  lacks the information'') actually holds on the data, empirically. Do
  not assume it.
\item
  \textbf{Hard-data, not curated-proxy.} Stress the claim on the
  genuinely-difficult regime it concerns (here: field degradation,
  occlusion, the expert-acknowledged-undecidable cases), scraped fresh,
  not on the curated benchmark, which is selected against difficulty.
\item
  \textbf{Symmetric verification.} If adversarial verification is used,
  give it a lens aimed at the negative: \emph{is this loss real, or an
  artifact of a weak setup?} That is not only \emph{is this win rigged?}
\item
  \textbf{Benchmark ≠ field.} State explicitly what the benchmark number
  does and does not bound. Separability on curated data bounds nothing
  about field safety, in either direction.
\end{enumerate}

\subsubsection{5.1 The negative-claim red-team (the lens our
verification
lacked)}\label{the-negative-claim-red-team-the-lens-our-verification-lacked}

Our adversarial-verification panels were effective against rigged
\emph{positives} and useless against an over-claimed \emph{negative},
because they never asked the negative-specific questions. We distill
those into a red-team a verification process should run against any
limitation/impossibility claim before it is trusted. Applied to our own
original paper, this is the protocol that would have refuted it in an
afternoon:

\begin{enumerate}
\def\labelenumi{\arabic{enumi}.}
\tightlist
\item
  \textbf{Strongest-baseline, not principled-baseline.} Name the single
  strongest method that could plausibly beat the claim; if the
  negative's baseline was chosen for principle or convenience (ours: a
  frozen, mean-pooled encoder for a clean proof), \emph{replace it with
  the strongest available and re-measure.} Bound the claim by the
  strongest baseline that still fails.
\item
  \textbf{Representation-invariance test.} Re-run the claim across ≥2
  paradigmatically different representations (here: masked-latent,
  self-distilled, contrastive). An \emph{information} limit must bind
  all of them; a \emph{representation} artifact will not. A
  within-instance gap between encoders refutes an information limit
  \emph{regardless of the absolute error level}. It is the single most
  diagnostic check, and the one we omitted.
\item
  \textbf{Antecedent audit.} If the negative invokes a theorem (``no
  function of X helps''), \emph{measure} its antecedent (``X lacks the
  discriminating information'') on the data; do not assume it. A valid
  theorem with an unverified antecedent is not evidence.
\item
  \textbf{Fresh hard-data, not curated proxy.} Test the claim on the
  genuinely-difficult regime it concerns, scraped fresh for the purpose,
  not the curated benchmark, which is selected against difficulty.
\item
  \textbf{No bare zeros.} Report every ``0 / no-method-helps'' cell with
  a one-sided upper confidence bound. A 0-of-\texttt{n} is not 0; under
  an asymmetric cost the difference between ``0\%'' and ``≤5\% with 95\%
  confidence'' can flip the conclusion.
\item
  \textbf{Cost-honest conclusion.} Re-state the conclusion against the
  \emph{upper} bound under the task's real loss: does the negative still
  hold if the true rate sits at the top of its confidence interval?
\end{enumerate}

The diagnostic asymmetry is item 2: a representation-invariance test is
inexpensive and conclusive, and was the one check the original
investigation never performed.

\begin{center}\rule{0.5\linewidth}{0.5pt}\end{center}

\subsection{6. What remains true (and the real safety
point)}\label{what-remains-true-and-the-real-safety-point}

The collapse of the information-limit thesis is not a safety
endorsement, and the corrected picture is itself a safety argument, just
a different one.

\begin{itemize}
\tightlist
\item
  \textbf{Strong encoders separate these species' \emph{distributions}
  very well; this does not make foraging safe.} Three concerns survive
  untouched: \emph{open-set true novelty} (a toxic species the model
  never saw), the \emph{cost asymmetry} (under \texttt{c\_PSC\ =\ 100},
  even a tiny tail false-eat can fail the decision-curve test), and the
  \emph{worst-case individual specimen} (the single \emph{Galerina}
  photographed at the moment it is truly identical to an
  \emph{Armillaria}, or an isolated button-stage \emph{Amanita} egg),
  which a random-sample 0\% cannot bound and which we could not obtain
  life-stage-filtered data to test. At full sampling the strong encoders
  cleared only a small fraction of held-out poisons (DINOv2 ≤0.7\% per
  cell, pooled 95\% UCB 0.4\%): low, but no longer a clean zero, and a
  random-sample rate this low still cannot bound the worst-case specimen
  above. Absence of a measured tail in a curated sample is not a field
  guarantee.
\item
  \textbf{The danger shifts from ``the model can't'' to ``the model is
  trusted because it can.''} A model that performs well even the
  microscope pairs invites exactly the over-trust the app-safety
  literature documents (higher confidence in AI \emph{reduces} correct
  decisions in high-risk identification). A near-perfect benchmark score
  is a \emph{reason for more caution}, not less, for a lethal decision.
\item
  \textbf{The formal propositions remain valid} as a statement about
  genuinely information-poor observations. They simply do not apply to
  the real lethal pairs we tested. They would apply to the worst-case
  corner above, which is the honest residual of the original thesis:
  narrow, real, and undemonstrated.
\end{itemize}

\begin{center}\rule{0.5\linewidth}{0.5pt}\end{center}

\subsection{7. Conclusion}\label{conclusion}

We built a careful, reproduced, adversarially-verified case that a
lethal recognition task is information-limited, and we were wrong: the
limit was a weak-baseline artifact. A stronger frozen encoder beat the
weak one \emph{within every pair}, which an information limit forbids,
across curated, field-degraded, and open-set regimes including the
textbook-undecidable pairs, and then again across six lethal look-alike
decisions in five biological groups, three of them re-examining
externally-published look-alike claims, two of visual
indistinguishability that a strong encoder violates and one of
systematic human misidentification that it instead corroborates. We
report the sequence because the failure is general: a limitation claim
is only as strong as the baseline that fails to beat it, curated
benchmarks understate field difficulty while overstating model limits,
and a verification discipline built to catch false positives will miss
false negatives. The contribution is methodological rather than
mycological: a worked demonstration, including its own collapse, that
safety-\emph{negative} claims demand the same adversarial stress as
safety-positive ones, with a checklist for supplying it. \emph{Research
vehicle, never a forager. When in doubt, DON'T EAT.}

\begin{center}\rule{0.5\linewidth}{0.5pt}\end{center}

\subsection{Appendix A: Formal statements (true, and empirically vacuous
for the real
pairs)}\label{appendix-a-formal-statements-true-and-empirically-vacuous-for-the-real-pairs}

These are the propositions the original investigation used to
\emph{explain} the limit. They are correct. The finding of §4(c) is that
their antecedent, ``the photograph lacks the discriminating
information'', is \textbf{false} for every real lethal pair we could
obtain (§3), so the propositions, though valid, bound nothing about
those pairs. We state them precisely both because they are the only
durable formal content and because seeing the antecedent explicitly
makes the vacuity clear.

\textbf{Setup.} Let \(X\) be the photograph on space \(\mathcal{X}\),
\(Y \in \{0,1\}\) the edibility label (\(0=\) edible, \(1=\) deadly),
and \(\varphi : \mathcal{X} \to \mathbb{R}^d\) a fixed feature map (the
encoder). Write the class-conditional feature laws
\(P_0 = \mathrm{Law}(\varphi(X)\mid Y{=}0)\) and
\(P_1 = \mathrm{Law}(\varphi(X)\mid Y{=}1)\), the posterior
\(\eta(z) = P(Y{=}1 \mid \varphi(X){=}z)\), and the prior
\(\pi = P(Y{=}1)\). A decision rule is a measurable
\(g : \mathbb{R}^d \to \{\text{eat},\text{abstain}\}\); the
\textbf{false-eat} rate is
\(\mathrm{FE}(g) = P(g{=}\text{eat}\mid Y{=}1)\) and the \textbf{yield}
is \(\mathrm{Yld}(g) = P(g{=}\text{eat}\mid Y{=}0)\). All three
statements concern a \emph{fixed} modality \(X\) (a single RGB
photograph).

\textbf{Proposition 1 (Coincidence / divergence floor).} Fix a target
yield \(\alpha\). Among rules with yield \(\ge \alpha\), the minimum
false-eat is the Bayes value \(\mathrm{FE}^*(\alpha)\), attained by
thresholding the likelihood ratio \(dP_1/dP_0\) (Neyman--Pearson). On
any region \(R\) where \(P_0 = P_1\): (i) every \(f\)-divergence and
every Wasserstein distance between \(P_0|_R\) and \(P_1|_R\) is \(0\);
(ii) \(dP_1/dP_0 = 1\), \(\eta \equiv \pi\), and
\(I(\varphi(X); Y \mid \varphi(X) \in R) = 0\); (iii) no rule separates
the classes on \(R\), so any rule clears poisons and edibles in equal
proportion there. \emph{Proof.} Each \(f\)-divergence and each \(W_p\)
vanishes iff the measures coincide; \(dP_1/dP_0 = 1\) gives
\(\eta = \pi\) by Bayes, and \(I = 0\) by conditional independence; with
\(\eta\) constant on \(R\) the Neyman--Pearson statistic is constant
there. \(\blacksquare\)

\textbf{Proposition 2 (Data-processing bound on post-hoc methods).} Let
\(Z = g(X)\) be \emph{any} measurable transform of the photograph: a
different representation, a generated or upscaled image, a topological
summary, a \emph{predicted} microscopy. Then \(Y \to X \to Z\) is a
Markov chain, so \(I(Z;Y) \le I(X;Y)\), and for every yield \(\alpha\),
\(\mathrm{FE}^*_Z(\alpha) \ge \mathrm{FE}^*_X(\alpha)\): no function of
the photograph beats the photograph's own Bayes floor. \emph{Proof.} The
data-processing inequality, plus monotonicity of the constrained Bayes
risk in the decision \(\sigma\)-algebra
(\(\sigma(Z) \subseteq \sigma(X)\)). \(\blacksquare\)

\textbf{Proposition 3 (When an extra channel helps: the independence
condition).} For corroboration rules (eat only when both \(X\) and an
auxiliary \(W\) clear), \(W\) lowers the false-eat below
\(\mathrm{FE}^*_X(\alpha)\) \textbf{iff} \(I(W; Y \mid X) > 0\). (i) A
redundant channel \(W = h(X, \xi)\) with \(\xi \perp\!\!\!\perp Y\) has
\(I(W;Y\mid X)=0\) and cannot help. (ii) An independent channel
\(W \perp\!\!\!\perp X \mid Y\) with per-class reliability \(\rho\)
multiplies the false-eat by \((1-\rho)\), at a yield cost. \emph{Proof.}
If \(W \perp\!\!\!\perp X \mid Y\), then \(\mathrm{FE}\) factors as
\(\mathrm{FE}^*_X(\alpha)\cdot(1-\rho)\); if \(I(W;Y\mid X)=0\), then
\(g_W\) is conditionally independent of \(Y\) given
\(\{g_X{=}\text{eat}\}\), leaving \(\mathrm{FE}\) unchanged.
\(\blacksquare\)

\textbf{The vacuity, stated plainly.} Propositions 1--3 are conditional:
\emph{if} a coincidence region of nonzero mass exists for a real lethal
pair, \emph{then} no method escapes it. Section 3 measured
\(\mathrm{FE}^*_X\) directly with a strong encoder and found it near
zero (\(\le 0.7\%\) per cell; pooled 95\% UCB \(0.4\%\)) across every
pair tested, including the textbook-undecidable ones. The coincidence
region of Proposition 1 therefore has at most negligible mass for these
pairs, and the precondition for Proposition 2 to bind (a large
\(\mathrm{FE}^*_X\)) does not hold. The theory is intact; it does not
describe the regime these pairs occupy. It would describe the worst-case
corner of Section 6 (a single genuinely-coincident specimen), which
remains undemonstrated for lack of life-stage-filtered data.

\subsection{Appendix B: The original
investigation}\label{appendix-b-the-original-investigation}

Preserved verbatim as the companion document (\texttt{MANUSCRIPT.md}):
the \textasciitilde18-method sweep, the decision-curve and conformal
apparatus, the VLM analysis, and the two self-caught rigged positives.
It is the specimen this paper dissects, not an endorsed result.

\subsection{Appendix C:
Reproducibility}\label{appendix-c-reproducibility}

All code and derived feature caches are provided as an anonymized,
view-only archive (contributor names hidden):
\texttt{osf.io/d3ft6/?view\_only=a7785622ad2e4e1995f2c830359447b3}. The
scripts named below are in its \texttt{code/} directory. The refutation:
\texttt{pairs\_encoders.py} (five pairs × three encoders, ±crop),
\texttt{hard\_regime.py} (field degradation + base occlusion + lookalike
open-set), \texttt{scrape\_hard.py} + \texttt{hard\_pairs.py} (the
needs-a-microscope pairs, ±substrate-strip crop). The two tables in §3
are the \textbf{full-image} re-run (\texttt{FULL=1}, which uses every
cached image and writes a separate \texttt{\_\_full} feature cache,
preserving the original 200/300-image arrays); both scripts print, per
cell, the number of unique held-out poison test specimens and a
one-sided 95\% Clopper-Pearson upper bound on the false-eat rate. The
second case study (§3.1): \texttt{scrape\_plants.py} (the deadly
Apiaceae) + \texttt{case2\_plants.py}. The §3.2 multi-taxon battery:
\texttt{scrape\_snakes.py} + \texttt{case3\_snakes.py} (coral-snake
mimicry), \texttt{scrape\_more.py} + \texttt{case4\_spiders.py}
(brown-recluse misID) and \texttt{case5\_frogs.py} (poison-frog
mimicry), \texttt{scrape\_berries.py} + \texttt{case6\_berries.py}
(deadly berries), each the same three-encoder protocol with
Clopper-Pearson bounds and a corrupt-image-robust loader. The
objective-vs-aggregation ablation (§4a): \texttt{ablation\_agg.py},
which re-encodes three hard pairs keeping patch tokens and compares
I-JEPA \{mean, max\} against DINOv2 \{CLS, mean, max\}. The original
investigation: the method-sweep scripts enumerated in the companion
document. All GBIF image scrapes are reproducible via the
\texttt{species/match} + \texttt{occurrence/search} endpoints in the
\texttt{scrape\_*.py} files.

The three frozen encoders are the public HuggingFace checkpoints
\texttt{facebook/ijepa\_vith14\_1k} (I-JEPA ViT-H/14, mean-pooled last
hidden state), \texttt{facebook/dinov2-large} (DINOv2 ViT-L/14, no
registers, CLS token), and \texttt{openai/clip-vit-large-patch14} (CLIP
ViT-L/14, pooled projection). Each image is preprocessed by that model's
own HuggingFace image processor, at its canonical resolution and
normalization. The decision probe of §3 is the MLP-ensemble risk net in
\texttt{model.py}, trained per seed on the 30\% split.

\textbf{AI use and verification.} The adversarial-verification panels
referred to throughout (the discipline of §2 and §4d, and the red-team
of §5.1) are panels of large-language-model agents run under a fixed
protocol: each finding is produced by one agent and then checked by
independent agents instructed to refute it, and only majority-surviving
findings are kept. We used large-language-model assistance for that
verification, for writing and editing, and for building the analysis and
scraping tooling. We did not use any model to generate, alter, or impute
data. All images are real GBIF photographs obtained through the
documented endpoints, and the synthetic-augmentation experiments are
described as experiments that did not remove the residual. Every
quantitative claim was produced by the named deterministic scripts, and
the Clopper-Pearson bounds that carry the conclusions were independently
re-derived and machine-checked. The authors verified each AI-assisted
claim against its artifact and take full responsibility for the content.

\subsubsection{Appendix C.1: Per-pair tables for the §3.2
battery}\label{appendix-c.1-per-pair-tables-for-the-3.2-battery}

The §3.2 summary collapses each case to a range; the full per-pair
results follow (10-seed AUROC / dangerous-cleared at 80\% edible-yield,
with one-sided 95\% Clopper-Pearson upper bound; \texttt{n} = unique
dangerous test specimens per seed, not seed-multiplied). The
``dangerous-cleared'' event is the venomous/toxic/model specimen cleared
as the harmless look-alike at the 80\%-yield threshold.

\textbf{Case 3: coral-snake mimicry (reptiles; \citep{pfennig2001}),
n=90:}

{\def\LTcaptype{none} % do not increment counter
\begin{longtable}[]{@{}
  >{\raggedright\arraybackslash}p{(\linewidth - 6\tabcolsep) * \real{0.2500}}
  >{\raggedright\arraybackslash}p{(\linewidth - 6\tabcolsep) * \real{0.2500}}
  >{\raggedright\arraybackslash}p{(\linewidth - 6\tabcolsep) * \real{0.2500}}
  >{\raggedright\arraybackslash}p{(\linewidth - 6\tabcolsep) * \real{0.2500}}@{}}
\toprule\noalign{}
\begin{minipage}[b]{\linewidth}\raggedright
pair
\end{minipage} & \begin{minipage}[b]{\linewidth}\raggedright
I-JEPA
\end{minipage} & \begin{minipage}[b]{\linewidth}\raggedright
DINOv2
\end{minipage} & \begin{minipage}[b]{\linewidth}\raggedright
CLIP
\end{minipage} \\
\midrule\noalign{}
\endhead
\bottomrule\noalign{}
\endlastfoot
coral / scarlet-kingsnake & 0.89 / 13.8\% (≤20.7\%) & 1.00 / 0.0\%
(≤3.3\%) & 0.98 / 0.7\% (≤5.2\%) \\
coral / milk-snake & 0.94 / 4.2\% (≤9.9\%) & 1.00 / 0.0\% (≤3.3\%) &
1.00 / 0.0\% (≤3.3\%) \\
coral / scarlet-snake & 0.84 / 22.4\% (≤30.6\%) & 1.00 / 0.0\% (≤3.3\%)
& 1.00 / 0.1\% (≤3.3\%) \\
\end{longtable}
}

\textbf{Case 4: brown-recluse misidentification (arachnids; Vetter),
n=90:}

{\def\LTcaptype{none} % do not increment counter
\begin{longtable}[]{@{}
  >{\raggedright\arraybackslash}p{(\linewidth - 6\tabcolsep) * \real{0.2500}}
  >{\raggedright\arraybackslash}p{(\linewidth - 6\tabcolsep) * \real{0.2500}}
  >{\raggedright\arraybackslash}p{(\linewidth - 6\tabcolsep) * \real{0.2500}}
  >{\raggedright\arraybackslash}p{(\linewidth - 6\tabcolsep) * \real{0.2500}}@{}}
\toprule\noalign{}
\begin{minipage}[b]{\linewidth}\raggedright
pair
\end{minipage} & \begin{minipage}[b]{\linewidth}\raggedright
I-JEPA
\end{minipage} & \begin{minipage}[b]{\linewidth}\raggedright
DINOv2
\end{minipage} & \begin{minipage}[b]{\linewidth}\raggedright
CLIP
\end{minipage} \\
\midrule\noalign{}
\endhead
\bottomrule\noalign{}
\endlastfoot
recluse / southern-house-spider & 0.73 / 48.9\% (≤58.0\%) & 0.98 / 2.0\%
(≤6.8\%) & 0.88 / 18.9\% (≤27.0\%) \\
recluse / wolf-spider & 0.93 / 10.3\% (≤16.8\%) & 1.00 / 0.0\% (≤3.3\%)
& 1.00 / 0.0\% (≤3.3\%) \\
recluse / hobo-spider & 0.84 / 26.0\% (≤34.2\%) & 1.00 / 0.0\% (≤3.3\%)
& 0.96 / 4.8\% (≤9.9\%) \\
\end{longtable}
}

\textbf{Case 5: poison-frog mimicry (amphibians; \citep{symula2001}
{[}Müllerian{]}; \citep{darst2006} {[}Batesian{]}):} the mimic/model
species have few GBIF photographs, so this case carries the smallest
test sets and widest bounds (n = 22--90). \emph{Allobates zaparo} was
topped up from 37 to 66 usable images by taking all media per occurrence
(GBIF has 37 occurrences but 71 distinct image URLs). Reported
direction: model / co-mimic cleared as the focal mimic.

{\def\LTcaptype{none} % do not increment counter
\begin{longtable}[]{@{}
  >{\raggedright\arraybackslash}p{(\linewidth - 8\tabcolsep) * \real{0.2000}}
  >{\raggedright\arraybackslash}p{(\linewidth - 8\tabcolsep) * \real{0.2000}}
  >{\raggedright\arraybackslash}p{(\linewidth - 8\tabcolsep) * \real{0.2000}}
  >{\raggedright\arraybackslash}p{(\linewidth - 8\tabcolsep) * \real{0.2000}}
  >{\raggedright\arraybackslash}p{(\linewidth - 8\tabcolsep) * \real{0.2000}}@{}}
\toprule\noalign{}
\begin{minipage}[b]{\linewidth}\raggedright
pair
\end{minipage} & \begin{minipage}[b]{\linewidth}\raggedright
n
\end{minipage} & \begin{minipage}[b]{\linewidth}\raggedright
I-JEPA
\end{minipage} & \begin{minipage}[b]{\linewidth}\raggedright
DINOv2
\end{minipage} & \begin{minipage}[b]{\linewidth}\raggedright
CLIP
\end{minipage} \\
\midrule\noalign{}
\endhead
\bottomrule\noalign{}
\endlastfoot
Müllerian: \emph{variabilis} / \emph{imitator} & 30 & 0.89 / 18.3\%
(≤35.7\%) & 0.96 / 4.7\% (≤14.9\%) & 0.93 / 10.3\% (≤23.9\%) \\
Müllerian: \emph{fantastica} / \emph{imitator} & 30 & 0.72 / 43.3\%
(≤59.8\%) & 0.99 / 0.3\% (≤9.5\%) & 0.98 / 2.7\% (≤14.9\%) \\
Batesian: \emph{Allobates} / \emph{Ameerega bilinguis} & 90 & 0.89 /
12.7\% (≤19.4\%) & 0.99 / 0.7\% (≤5.2\%) & 0.98 / 0.9\% (≤5.2\%) \\
Batesian: \emph{Allobates} / \emph{Ameerega parvula} & 22 & 0.82 /
33.6\% (≤51.5\%) & 0.93 / 10.0\% (≤25.9\%) & 0.92 / 16.4\% (≤36.9\%) \\
\end{longtable}
}

\textbf{Case 6: deadly berries (plant fruit; foraging literature),
n=90:}

{\def\LTcaptype{none} % do not increment counter
\begin{longtable}[]{@{}
  >{\raggedright\arraybackslash}p{(\linewidth - 6\tabcolsep) * \real{0.2500}}
  >{\raggedright\arraybackslash}p{(\linewidth - 6\tabcolsep) * \real{0.2500}}
  >{\raggedright\arraybackslash}p{(\linewidth - 6\tabcolsep) * \real{0.2500}}
  >{\raggedright\arraybackslash}p{(\linewidth - 6\tabcolsep) * \real{0.2500}}@{}}
\toprule\noalign{}
\begin{minipage}[b]{\linewidth}\raggedright
pair
\end{minipage} & \begin{minipage}[b]{\linewidth}\raggedright
I-JEPA
\end{minipage} & \begin{minipage}[b]{\linewidth}\raggedright
DINOv2
\end{minipage} & \begin{minipage}[b]{\linewidth}\raggedright
CLIP
\end{minipage} \\
\midrule\noalign{}
\endhead
\bottomrule\noalign{}
\endlastfoot
elderberry / pokeweed & 0.88 / 19.9\% (≤28.2\%) & 1.00 / 0.0\% (≤3.3\%)
& 0.99 / 0.0\% (≤3.3\%) \\
wild-grape / moonseed & 0.91 / 10.8\% (≤18.1\%) & 1.00 / 0.0\% (≤3.3\%)
& 1.00 / 0.3\% (≤3.3\%) \\
bilberry / deadly-nightshade & 1.00 / 0.2\% (≤3.3\%) & 1.00 / 0.0\%
(≤3.3\%) & 1.00 / 0.0\% (≤3.3\%) \\
\end{longtable}
}

At n=90 a ``0.0\%'' cell upper-bounds at 3.3\%; on the small frog pairs
a zero cell bounds only at \textasciitilde9.5\% (n=30) to
\textasciitilde12.7\% (n=22), and DINOv2 itself reaches 10.0\% (≤25.9\%)
on the n=22 \emph{parvula} pair. The frog case follows the same
weak→strong ordering in every pair (DINOv2 below I-JEPA point-wise
throughout), but its bounds are wide enough that only the n=90
\emph{bilinguis} pair is statistically decisive: read the frog evidence
as suggestive, not confirmatory. The other four groups (n=90--150) are
firm.

\begin{center}\rule{0.5\linewidth}{0.5pt}\end{center}
